% Reconstructed text from V5 PDF page 2; non-native source
\par\smallskip\noindent\rule{\textwidth}{0.4pt}\par{\sffamily\footnotesize Original V5 page 2 -- recovered text}\par\begin{flushleft}\footnotesize\ttfamily
\par\smallskip
\par\smallskip
\noindent Diagnostics and Serving Infrastructure for MAS 2\par
\par\smallskip
\par\smallskip
\par\smallskip
\noindent 1 Introduction\par
\noindent Foundation model-based multi-agent systems (MAS), in which multiple Large Language Model\par
\noindent (LLM) agents collaborate to solve complex tasks, are increasingly moving from demonstrations\par
\noindent toward deployed workflows. As systems such as MetaGPT [13], CAMEL [15], ChatDev [30],\par
\noindent and Internet of Agents [7] operate within shared environments, reliability and efficiency de-\par
\noindent pend on operational layers that are less mature than orchestration itself. As these systems be-\par
\noindent come architecturally sophisticated, reliability and efficiency depend on operational layers that\par
\noindent lag behind the orchestration layer: the diagnostics layer, which supports failure localization,\par
\noindent observability, and safety intervention; and the infrastructure layer, which manages memory,\par
\noindent scheduling, and resources for multi-agent workloads.\par
\noindent  Recent multi-agent applications span collaborative coding, scientific discovery, financial\par
\noindent analysis, and social simulation. Within these applications, failure patterns are qualitatively\par
\noindent different from single-agent errors: MAST [5] identifies 14 recurring failure modes in MAS,\par
\noindent motivating runtime diagnostics rather than evaluation only at the final task output. When a\par
\noindent multi-agent workflow fails, the central question is not just what failed but which agent caused\par
\noindent the failure, at which step, and whether the serving infrastructure’s resource management deci-\par
\noindent sions contributed to or masked the failure. Recent preprints illustrate adjacent model-level prob-\par
\noindent lems that may arise in agent-enabled workloads: finite-beam reachability and continuous-query\par
\noindent generation in generative recommendation, attribution-based forgetting in LoRA adaptation, and\par
\noindent reading and reconciling temporal evidence in time-series forecasting. These are contextual ex-\par
\noindent amples rather than multi-agent diagnostic or serving systems, and are not included in our core\par
\noindent system tables.\par
\par\smallskip
\noindent 1.1 Scope and Positioning\par
\noindent Existing surveys largely treat diagnostics and infrastructure separately. General agent surveys\par
\noindent cover architectural patterns and collaboration mechanisms; AgentOps addresses telemetry and\par
\noindent operational practice; offline agent evaluation benchmarks focus on task quality; single-model\par
\noindent serving surveys address latency and throughput optimization; and MAS collaboration/failure\par
\noindent surveys examine interaction patterns and failure modes. In our curated corpus, we review 48\par
\noindent entries across four categories (12 failure-localization, 6 observability, 8 safety, 22 infrastructure)\par
\noindent and 12 evaluation benchmarks (60 table entries; Appendix A). The resulting literature exposes\par
\noindent complementary pieces of the runtime stack, but the relationship between the state required for\par
\noindent diagnosis and the state managed by serving policies remains under-analyzed.\par
\noindent  This review examines the interface between these two selected runtime functions. The title\par
\noindent does not imply exhaustive treatment of orchestration, storage, networking, security operations,\par
\noindent or all layers of a production stack. Our analytical contribution is a state-dependency and cost\par
\noindent framework, not a claim that the reviewed mechanisms themselves are new. Table 1 contrasts\par
\noindent our scope against five adjacent review areas, specifying the conceptual dimension this review\par
\noindent adds to each.\par
\noindent  Three broad system properties recur across the literature reviewed here: (1) distributed\par
\noindent execution across multiple agents, which creates traces spanning agent boundaries; (2) inter-\par
\noindent agent communication, which allows failures to propagate through communication channels;\par
\noindent and (3) inter-agent memory sharing, where prompts, tools, and cached state may be reused\par
\noindent across the system. Section 5 derives serving-side implications from these properties, high-\par
\noindent lighting sparse activation, invocation distance, and inter-agent memory sharing as the workload\par
\noindent characteristics that directly affect memory placement, scheduling, and cache policy.\par
\end{flushleft}
